\section{From puncture data to convergence of curves}
\label{sec:identification}

Proposition~\ref{prop:tightness} retains every macroscopic loop in a subsequential limit, but it allows that loop to touch the boundary or retrace part of its path. We now exclude these possibilities. Theorem~\ref{thm:cylindrical} first assigns a canonical $\CLE_4$ loop to every limiting loop visible from the punctures. Proposition~\ref{prop:local-opposite} then rules out all motion away from the assigned trace, as well as backtracking along that trace.

Throughout this section, diameters and uniform convergence are measured after applying $F$. A path in the closed disk is equivalently a path in
\[
 \overline{\HH}^{\,\infty}:=\overline{\HH}\cup\{\infty\}
\]
with distance $|F(z)-F(w)|$. Coordinate lines and intersection indices will only be used in bounded parts of the half-plane. Write $\mathbb T=\mathbb R/\mathbb Z$ for the parameter circle.

\subsection{A joint limit and its puncture signatures}

Fix any sequence $\delta_n\downarrow0$. By Proposition~\ref{prop:tightness}, pass to a subsequence with convergent ordered, parametrized loop collections. On a suitable probability space their representatives satisfy
\begin{equation}\label{eq:representative-convergence}
 F\circ\gamma_{n,j}\longrightarrow F\circ\gamma_j
 \quad\text{uniformly on }\mathbb T
\end{equation}
for every retained nonconstant entry $j$. The convergence also retains every loop above each positive diameter, with multiplicity. One may pad a finite list by constants and then discard constant limiting entries. In particular, for each $a>0$ all loops of diameter at least $a$ lie in a finite initial portion of the lists, eventually. Indeed, choose an index whose limiting diameter is less than $a/2$; its approximating diameter is eventually less than $a$, and diameter ordering gives the same bound for every later entry. We may enlarge this joint realization by further tight coordinates and take further subsequences below.

We choose the test lines before choosing the punctures. For a continuous real function on an interval, the set of its local maximum and minimum \emph{values} is countable. Indeed, each such value is the maximum or minimum on some compact interval with rational endpoints inside a corresponding extremum neighborhood. This argument includes nonstrict extrema and flat intervals. Apply it to the coordinates of all $\gamma_j$ in countably many finite parameter charts. Fubini's theorem gives deterministic countable dense families of vertical lines and positive-height horizontal lines such that, almost surely, none of their levels is a local extremal value of any $\gamma_j$. We also avoid every lattice coordinate in the chosen mesh sequence. Consequently, at each visit of a limiting path to a test line, every parameter neighborhood contains points on both strict sides of that line.

On every test line fix successively finer uniform grids, with mesh tending to zero, and bounded testing ranges increasing to cover its finite part. Choose grid origins so that their endpoints avoid the graphs for every $n$. For vertical lines include the intervals meeting height zero; endpoints at nonpositive height have inside indicator zero for every discrete loop. Let
\[
 Q=\{q_1,q_2,\ldots\}\subset\HH
\]
contain all interior grid endpoints and be dense in $\HH$, with all added points chosen off the countable union of the lattice graphs. All these choices are deterministic. For each fixed $q\in Q$ and all sufficiently small meshes, $q$ lies in a unique bounded lattice face. Replacing it by that face center preserves every loop membership, and the centers converge to $q$, as required in Theorem~\ref{thm:cylindrical}. By Proposition~\ref{prop:cle-facts}, almost surely no point of $Q$ lies on a canonical $\CLE_4$ trace.

For each discrete loop define
\[
 a_{n,j}(q)=\1\{q\text{ lies inside }\gamma_{n,j}\},\qquad q\in Q.
\]
Set all these indicators to zero for a padded constant entry.
Adjoin these binary coordinates and the finite-puncture multiset records to the joint subsequence. The binary product is compact; the finite-puncture coordinates are tight by Theorem~\ref{thm:cylindrical}. The ambient space is the Polish product
\[
 C(\mathbb T,\overline{\DD})^{\mathbb N}
 \times\{0,1\}^{\mathbb N\times\mathbb N}
 \times\prod_{m\ge1}\mathcal R_m,
\]
where $\mathcal R_m$ is the countable discrete space of finite reduced multiset records for $q_1,\ldots,q_m$. Choose a compact restriction in each coordinate with a summable allocation of the error probability. Their product is compact, and the union bound proves joint tightness without any independence assumption. The Prokhorov and Skorokhod theorems \cite[Theorems~4.1 and~3.3]{Billingsley} therefore supply a further subsequence and a realization on one probability space such that
\begin{equation}\label{eq:signature-stability}
 a_{n,j}(q)=a_j(q)\quad\text{eventually, for each fixed }j,q.
\end{equation}
The limiting path marginal is unchanged by this further extraction and realization, so the probability-one generic-line property chosen above still holds. We call $\{q\in Q:a_j(q)=1\}$ the \emph{signature} of the entry $j$. This definition does not assert that the signature is already realized by the winding of $\gamma_j$.

\begin{lemma}[Identification of the visible signatures]\label{lem:visible-signatures}
There is a coupling with a standard nested $\CLE_4$ collection $\mathcal C$ such that the limiting entries whose signatures contain at least two points are in bijection with the loops of $\mathcal C$, preserving all memberships in $Q$. Every other nonconstant entry has empty signature.
\end{lemma}

\begin{proof}
A loop enclosing two distinct points $q,q'$ has disk diameter at least $|F(q)-F(q')|$: its enclosed disk domain lies in the convex hull of its boundary. Thus all entries containing this pair belong to an eventually finite initial portion of the ordered lists. Their stabilized memberships in every larger finite puncture set agree with the corresponding limiting multiset record. Theorem~\ref{thm:cylindrical}, applied to the largest of any finite family of puncture sets, gives the joint law of all these records.

These finite records determine the full collection of signatures containing at least two points. Fix a pair $q,q'$. There are finitely many signatures containing it. Their restrictions to increasing finite subsets of $Q$ determine that finite multiset of binary sequences: one can choose consistent identifications by passing through the finitely many permutations of its entries. To combine the records without counting a signature again for each pair it contains, enumerate the unordered pairs in $Q$ and assign each full signature to its first contained pair. Its multiplicity is the one recovered from that pair's record. This determines the entire signature multiset, which therefore has the same law as the corresponding multiset for nested $\CLE_4$.

To couple an actual collection $\mathcal C$ with these records, realize standard nested $\CLE_4$ on a standard Borel probability space, for example using its loop-soup construction and countably many independent copies for the nested generations. Condition that underlying randomness on its countable puncture record, using the regular conditional distribution theorem \cite[Theorems~4.1.17--4.1.18]{Durrett}, and sample the resulting conditional law at our limiting record. More explicitly, if $Y$ denotes all the already coupled paths and records, $R(Y)$ their limiting puncture record, and $K(r,d\omega)$ this conditional kernel for the canonical construction, extend the probability law by
\[
 \PP_Y(dy)\,K(R(y),d\omega).
\]
Integrating out $\omega$ preserves the entire original joint law of $Y$, while integrating over $y$ gives the standard canonical law because $R(Y)$ has its required marginal distribution. The two puncture records agree almost surely by the defining property of the conditional kernel. This gives the required coupling without any assumption about completeness of the curve quotient.

Every canonical loop has a nonempty open interior and therefore contains two points of $Q$. Distinct canonical loops have distinct signatures, because their interiors differ on an open set. This gives the claimed bijection for signatures of size at least two, with no duplicate visible entries.

It remains to exclude a singleton signature. Suppose a nonconstant additional entry has signature $\{q\}$ and diameter $a>0$. Proposition~\ref{prop:cle-facts} supplies infinitely many canonical loops around $q$. Every such loop contains another sampled point $q'$, so its corresponding discrete loop eventually contains $q$ and $q'$, whereas the additional discrete loop contains $q$ but not $q'$. Discrete loops are disjoint or nested. The canonical loop's discrete representative must therefore enclose the additional one. Its limiting diameter is at least $a$. Applying this argument to every finite number of the infinitely many canonical loops contradicts the finiteness of the limiting collection above diameter $a/2$. A nonvisible nonconstant entry consequently has empty signature.
\end{proof}

\subsection{Fixed tests followed by a slowly refining grid}

We next combine signature stabilization with the geometric estimate. The order of limits matters: no puncture convergence theorem with colliding or mesh-dependent punctures will be used.

Enumerate the test lines and their fixed bounded testing ranges, bounded windows with rational coordinate bounds, and positive rational side buffers. For each finite initial batch, Proposition~\ref{prop:local-opposite} says that the probability of any opposite-sign violation on its grid intervals tends to zero when the grid mesh tends to zero \emph{after} the lattice-mesh upper limit. Choose a sufficiently fine grid level for the first batch, then a lattice-mesh threshold making this failure probability small. At that fixed grid level, there are only finitely many endpoint indicators for any fixed finite prefix of the loop list. Equation~\eqref{eq:signature-stability} lets us decrease the lattice-mesh threshold so that failure of their simultaneous stabilization is also small. Repeat with increasing batches and loop prefixes, using a summable sequence of error probabilities.

Passing to the resulting subsequence and applying Borel--Cantelli gives grid levels $r_n\to\infty$ with the following almost-sure property. For every fixed finite batch, eventually its level-$r_n$ intervals have no local opposite-sign violation, and all indicators used for each fixed finite prefix of loops agree with their stabilized signatures. The grids themselves were fixed before any of these lattice-mesh thresholds were chosen. We shall use this procedure again after identifying the traces, adding further countably many coordinates which are then known to stabilize.

Here and below, a directed \emph{subpath} means a restriction to a parameter interval; it need not be injective in the limit. For a directed bounded path $\alpha$ and an oriented line interval $I=[u,v]$, with $u,v$ off its image and the path endpoints off the line, denote its oriented intersection index by $\iota_I(\alpha)$, with the convention of Proposition~\ref{prop:local-opposite}. Orient $I$ consistently with that convention. For polygonal paths this is the signed crossing sum. In general, close $\alpha$ by a path avoiding $I$ and set
\[
 \iota_I(\alpha)=\operatorname{wind}(\alpha\text{ with its closing path},u)
             -\operatorname{wind}(\alpha\text{ with its closing path},v).
\]
The value does not depend on the closing path, since a closed path avoiding $I$ has equal winding about its endpoints. This definition is additive under concatenation and stable under uniform perturbations avoiding $u,v$, with endpoints staying off the test line. For a Jordan loop it is, up to the loop's orientation, the difference of the two inside indicators. In particular, it belongs to $\{0,1,-1\}$.

The following elementary partition observation permits use of Proposition~\ref{prop:local-opposite}, despite potentially infinitely many line hits.

\begin{lemma}[Partition into eligible portions]\label{lem:opposite-partition}
Let $\gamma:\mathbb T\to\overline{\HH}^{\,\infty}$ be continuous, and let $L$ be a coordinate line whose level is not a local extremal value of the corresponding coordinate of $\gamma$ at finite locations. For every bounded segment $K\subset L$, there is a finite parameter partition such that every piece meeting $K$ has endpoints on opposite strict sides of $L$ and lies in a bounded window. The other pieces stay off $K$ under sufficiently small uniform perturbations. The same assertion holds while retaining a prescribed bounded subpath with opposite-side endpoints as one unsplit piece, provided that subpath is considered separately when choosing the small-oscillation partition.
\end{lemma}

\begin{proof}
Take a finite small-oscillation partition, so fine in the disk metric that pieces meeting a neighborhood of $K$ lie in a bounded window. Choose their endpoints off $L$; such parameter values are dense by the assumed absence of extremal levels. A piece touching $L$ with endpoints on the same side contains a point on the other side, again by that assumption. Split it once at this point. The two resulting pieces have opposite-side endpoints. Pieces already having opposite-side endpoints need no alteration. Each remaining piece has compact image disjoint from $K$, and hence stays disjoint under sufficiently small perturbations. Enlarging $K$ slightly at the start avoids issues at its ends. To retain a prescribed subpath, begin with its endpoints in the partition and perform this construction on the complementary parameter interval. The relevant bounded pieces have finitely many endpoints, so their strict distances from $L$ have a positive minimum.
\end{proof}

Thus all pieces relevant to a fixed bounded test range satisfy some fixed positive side buffers. On sufficiently close discrete approximations they satisfy the same conditions. Realized partitions do not require an uncountable extension of the probability estimate: that estimate already quantifies over \emph{all} subpaths within the fixed window and buffers. The countable families of windows and rational buffers cover every partition just constructed.

\subsection{Excluding invisible paths and identifying the traces}

Fix a nonconstant limiting entry $\gamma_j$. For a bounded testing range on a selected line, apply Lemma~\ref{lem:opposite-partition}. On every interval of the slowly refining grids, all nonzero indices of the resulting discrete portions have the same sign, by Proposition~\ref{prop:local-opposite}. Their sum is the whole-loop index. Consequently,
\begin{equation}\label{eq:no-cancellation}
 \iota_I(\alpha_n)\ne0
 \quad\Longrightarrow\quad
 a_{n,j}(u)\ne a_{n,j}(v)
 \qquad(I=[u,v])
\end{equation}
for every retained portion $\alpha_n$ of this partition. Endpoints outside $\HH$ are understood to have indicator zero.

There is a useful local consequence. Suppose a subpath of $\gamma_j$ stays in a small finite ball and crosses a selected line from one strict side to the other. Retain that subpath in the partition. Its discrete approximation has net signed crossing $+1$ or $-1$ with the entire line. Since the approximation stays in a slightly enlarged ball, that crossing is the sum of its indices on grid intervals in the enlarged ball. At least one of those indices is nonzero. Equation~\eqref{eq:no-cancellation} then forces an endpoint-indicator switch on one such interval. For sufficiently fine grids all relevant intervals lie in any fixed larger ball. By the simultaneous stabilization already arranged, the same switch occurs in the limiting signature.

\begin{lemma}[Identification of traces and degree]\label{lem:trace-identification}
There is no nonconstant entry with empty signature. Each remaining limiting path has exactly the trace of its corresponding loop of $\mathcal C$ and traverses it with absolute degree one.
\end{lemma}

\begin{proof}
An empty signature permits no endpoint-indicator switch at any interior grid endpoints, nor at exterior endpoints, whose indicator is zero. The local consequence above therefore excludes every finite subpath crossing a test line. Every genuine finite movement can be separated by one of the dense coordinate lines. A nonconstant path visiting infinity also has a nonconstant finite portion. Thus an empty-signature path would have to be constant, a contradiction.

Now let $C\in\mathcal C$ be the canonical loop paired with $\gamma_j$. In any finite open ball disjoint from $C$, the canonical inside indicator is constant. If the ball meets the real boundary, that value is zero, including at exterior points. Hence the signature allows no switch on sufficiently interior grid intervals in this ball. Any motion of $\gamma_j$ in a smaller ball would give the forbidden switch just proved. A nonconstant path visiting a point off $C$ must make such a motion: even if it pauses at that point, continuity forces motion in the complementary open set when it reaches or leaves the pause. A visit to infinity similarly has finite portions off the compact curve $C$. We conclude that the trace of $\gamma_j$ is contained in $C$; in particular it is finite and interior.

Choose $q\in Q$ inside $C$. It is off $C$, and its stabilized indicator is one. Uniform convergence now takes place in a bounded Euclidean neighborhood of $C$, and preserves winding about $q$. Every sufficiently late simple approximating loop has absolute winding one about $q$, so the limiting path has absolute winding one. Identifying $C$ with a circle shows that $\gamma_j$ has degree $+1$ or $-1$ as a map into $C$. A nonzero-degree map of the circle is surjective. Its trace is therefore all of $C$.
\end{proof}

This has identified the sets traced by all macroscopic loops and their multiplicities. It has not yet identified their classes in the curve metric: a degree-one traversal can move backward over a substantial arc before continuing forward. The final geometric estimate excludes precisely this behavior.

\subsection{The exclusion of backtracking}

We first make explicit a local fact about Jordan curves. No smoothness or transverse derivative is assumed.

\begin{lemma}[A generic coordinate line sees both Jordan domains]\label{lem:jordan-line}
Let $C$ be a Jordan curve, let $p\in C$, and let $L$ be a vertical or horizontal line through $p$. Suppose every open Jordan subarc around $p$ contains points on both strict sides of $L$. Then every neighborhood of $p$ along $L$ contains points in both complementary domains of $C$.
\end{lemma}

\begin{proof}
By the Schoenflies theorem \cite[Section~1]{Cairns}, choose a plane homeomorphism $h$ taking the unit circle to $C$, with $h(q)=p$. Given $\eps>0$, take a small round disk $B$ centered at $q$ such that $h(B)\subset B_\eps(p)$. The images of
\[
 B\cap\DD\quad\text{and}\quad
 B\cap(\mathbb C\setminus\overline\DD)
\]
are connected, and each has the open subarc $h(B\cap\partial\DD)$ in its closure. If one complementary domain of $C$ missed $L\cap B_\eps(p)$, its corresponding connected image would avoid $L$ and hence lie in one strict half-plane. Its closure would put that entire Jordan subarc in the corresponding closed half-plane, contrary to the hypothesis. Both complementary domains therefore meet $L\cap B_\eps(p)$.
\end{proof}

All fixed grid endpoints now avoid the limiting traces: interior endpoints belong to $Q$, and the traces themselves lie in $\HH$. Fix a finite grid, a finite prefix of retained nonconstant entries, and finitely many directed parameter intervals with rational endpoints. For each such interval whose two limiting path endpoints are strictly off its test line, uniform convergence preserves its index on every interval of the grid for all sufficiently large $n$. A parameter interval with a limiting endpoint on the line imposes no test. There are only finitely many tests, and each included pair has a positive realized distance from its line. Their simultaneous failure probability therefore tends to zero by bounded convergence.

Repeat the earlier diagonal selection with these additional index equalities. At stage $k$, choose a finer deterministic grid for the first $k$ geometric tests, then keep that grid fixed while choosing $n_k$ for the first $k$ retained entries and rational parameter intervals. The geometric violations, membership mismatches, and new index mismatches can together be given probability less than $2^{-k}$. Borel--Cantelli yields their eventual exclusion for every fixed test, entry, and rational interval whose limiting endpoints are off the line. The indices can be added only at this stage: trace identification has supplied the avoidance of grid endpoints needed for their stability. This second extraction preserves all the earlier almost-sure conclusions and again fixes each finite grid before taking the lattice mesh sufficiently small.

\begin{lemma}[No reverse passage]\label{lem:no-backtracking}
Each path $\gamma_j$ is a weakly monotone once-around traversal of its canonical Jordan curve. In particular, it has zero curve distance from a simple parametrization of that curve.
\end{lemma}

\begin{proof}
Reverse the orientation if needed so that the degree is $+1$. Let $\eta:\mathbb R/\mathbb Z\to C$ be a positively oriented homeomorphism. There is a continuous lift $\phi:\mathbb R\to\mathbb R$ with
\[
 \gamma_j(t)=\eta(\phi(t)\bmod1),\qquad
 \phi(t+1)=\phi(t)+1.
\]
If $\phi$ is not nondecreasing, choose two lift values $\alpha<\beta$, with $\beta-\alpha<1$, that it crosses in the decreasing direction. The points $\eta(\alpha\bmod1)$ and $\eta(\beta\bmod1)$ are distinct, so one planar coordinate separates them. Select a test line $L$ strictly between those coordinate values.

First and last hitting times extract a downward passage on $[s,t]$ from $\beta$ to $\alpha$ whose lift remains in $[\alpha,\beta]$. Since $\phi(s+1)=\beta+1$, necessarily $t<s+1$. The same identity guarantees a subsequent upward passage from $\alpha$ to $\beta$ before time $s+1$; extract it in the same way. Thus both passages lie in one period and have disjoint time interiors, their endpoints are on opposite sides of $L$, and both traces equal the proper Jordan subarc
\[
 K=\eta([\alpha,\beta]).
\]
As paths within $K$, they are homotopic relative to endpoints to its two opposite traversals, by linear homotopy of their lifts.

Choose a point $p\in\eta((\alpha,\beta))\cap L$. It is not a local coordinate extremum on $C$: such an intrinsic extremum would give a local extremum of the same coordinate at any time when $\gamma_j$ visits it, contrary to the choice of the test line. Take a disk around $p$ meeting $C$ only in the common open subarc $\eta((\alpha,\beta))$. Lemma~\ref{lem:jordan-line} supplies open intervals on $L$ within this disk lying in the two different Jordan domains. At every sufficiently fine grid level there are grid endpoints in both intervals. Among consecutive endpoints between them, some adjacent pair $u,v$ has different inside indicators. Its interval $I=[u,v]$ remains in the disk, and all intersections of $C$ with $I$ lie in the common open subarc.

The two passages have opposite nonzero indices on $I$. Indeed, the homotopies within $K$ identify their indices with opposite traversals of $K$, and the rest of the Jordan curve avoids $I$. Additivity thus identifies either index, up to sign, with the full Jordan-curve index on $I$, which is the difference of its endpoint indicators and has absolute value one.

We may replace the four passage endpoint times by nearby rational times, perturbing inward so that the time interiors remain disjoint. Choose these perturbations so that their removed path pieces remain near the original endpoints, away from the disk just selected. Their endpoints still lie on opposite strict sides of $L$, and their indices on every interval in that disk are unchanged. Both resulting portions lie in a fixed bounded window and have fixed positive side buffers. The repeated slow-grid extraction transfers their opposite nonzero indices to the discrete loop for all sufficiently late members of the subsequence. This contradicts Proposition~\ref{prop:local-opposite}.

Thus $\phi$ is nondecreasing. A nondecreasing degree-one parametrization differs from a homeomorphic parametrization only by pauses in the curve metric. Explicitly, the strictly increasing lifts
\[
 \phi_\eps(t)=\frac{\phi(t)+\eps t}{1+\eps}
\]
converge uniformly on one period to $\phi$ and satisfy $\phi_\eps(t+1)=\phi_\eps(t)+1$. Their circle homeomorphisms give zero distance to the simple parametrization $\eta$.
\end{proof}

\subsection{The required collection space and the full limit}

Every nonconstant limiting entry is now one canonical $\CLE_4$ loop as a curve, and every canonical loop occurs once. Proposition~\ref{prop:tightness} controls the positive-diameter tails, so each closed-disk limit has exactly the nested $\CLE_4$ law. Since the starting mesh sequence was arbitrary, the subsequence criterion gives full convergence in this ambient space.

The remaining task is to obtain compact containment in the interior collection space. The following lemma separates this topological step from the identification of the curves. Write $X$ for the closed-disk collection space with metric $\rho$, and $S\subset X$ for its subspace of interior collections.

\begin{lemma}[Restriction to interior collections]\label{lem:interior-restriction}
Let $\mu_\delta$, $0<\delta\le1$, and $\mu$ be probability laws on $X$ supported on $S$. Suppose that $\mu_\delta\Rightarrow\mu$ in $X$ as $\delta\downarrow0$, that the family $(\mu_\delta)_{0<\delta\le1}$ is uniformly tight in $X$, and that $(\mu_\delta)_{a\le\delta\le1}$ is uniformly tight in $S$ for every $a>0$. Then $\mu_\delta\Rightarrow\mu$ in $S$, and the entire family is uniformly tight in $S$.
\end{lemma}

\begin{proof}
For $m\ge1$, let
\[
 B_m=\{\mathcal L:\text{some loop in }\mathcal L
       \text{ has diameter at least }1/m
       \text{ and meets }\partial\DD\}.
\]
Each $B_m$ is closed. To see this, along convergent partial matchings a witnessing loop must match one of the finitely many limiting loops of diameter at least $1/(2m)$. Pass to a subsequence using the same limiting entry. Uniform curve convergence preserves the diameter bound and boundary contact. The space $S$ of interior loop collections is consequently
\begin{equation}\label{eq:interior-space}
 S=\bigcap_{m\ge1}B_m^{\,c}.
\end{equation}
An open subset of $S$ is $G\cap S$ for an ambient open set $G$. Since all the laws give $S$ probability one, the open-set Portmanteau inequality in $X$ gives the same inequality in $S$, proving the asserted weak convergence.

Fix $\eps>0$ and choose an ambient compact set $K$ with $\mu_\delta(K)\ge1-\eps/2$ for every $0<\delta\le1$. For each $m$, a $\mu$-distributed collection has strictly positive distance from $B_m$. Ambient weak convergence, applied to a sufficiently small closed distance sublevel set, therefore gives a radius $a_m>0$ and a threshold $\delta_m>0$ such that
\[
 \mu_\delta\{\mathcal L:\dist(\mathcal L,B_m)<a_m\}
     <\eps\,2^{-m-1},\qquad 0<\delta<\delta_m.
\]
For $\delta_m\le\delta\le1$, uniform tightness in $S$ supplies a compact subset of $S$ carrying probability greater than $1-\eps2^{-m-1}$ for every such law. This compact set has positive distance from the closed set $B_m$. Decreasing $a_m$ accordingly makes the displayed bound hold for every $0<\delta\le1$. Hence
\[
 K\cap\bigcap_{m\ge1}
       \{\mathcal L:\dist(\mathcal L,B_m)\ge a_m\}
\]
is compact, lies in $S$, and has probability at least $1-\eps$ under every $\mu_\delta$. This proves uniform tightness in $S$.
\end{proof}

For our laws, only tightness on compact positive mesh intervals remains to be checked. On any fixed lattice the possible finite cycles form a countable set, and each cycle can occur at most once. At each positive diameter cutoff, choose finitely many such cycles to capture every loop above that cutoff with arbitrarily high probability. Make the errors summable over decreasing cutoffs. The collections of distinct fixed-lattice cycles satisfying all these restrictions form a compact subset of $S$: diagonalize membership of the possible cycles, and at any fixed cutoff match the finitely many retained cycles exactly. The unmatched tails have vanishing diameter. Thus the fixed lattice law is tight in $S$.

Moreover, $\mathcal L_\delta$ has exactly the law of $\delta\mathcal L_1$. The maps
\[
 T_\delta=F\circ(z\mapsto\delta z)\circ F^{-1}
\]
are jointly continuous on $[a,b]\times\overline\DD$ for every $0<a<b<\infty$. They induce jointly continuous maps on collections as well: their common modulus of continuity controls the distances of matched loops and the diameters of unmatched loops, while a change in $\delta$ changes $T_\delta$ uniformly. Images of a compact restriction for $F(\mathcal L_1)$ under these maps are compact subsets of $S$. The laws with $\delta\in[a,b]$ are therefore uniformly tight in $S$.

\begin{proof}[Completion of the proof of Theorem~\ref{thm:main}]
We have identified every subsequential closed-disk limit as standard nested $\CLE_4$, with every macroscopic loop matched and with multiplicity one, and hence obtained full ambient convergence. Proposition~\ref{prop:tightness} supplies uniform ambient tightness for $0<\delta\le1$. The discrete laws and the identified limit are supported on $S$, and the preceding scaling argument gives tightness in $S$ on each compact positive mesh interval. Lemma~\ref{lem:interior-restriction} therefore gives the claimed convergence and tightness in the interior collection topology.
\end{proof}
